\section{Introduction}
\label{sec:intro}

Low-cost, open-source potentiostats have made electroanalytical measurement
accessible far beyond the benchtop, in point-of-care, field and teaching
settings~\cite{rowe2011cheapstat, jenkins2019abe, brown2022acestat,
bautista2025helpstat, anshori2025ganestat}. Integrated electrochemical analog
front ends, in particular the Analog Devices AD5940/AD5941, now place a
complete potentiostatic control loop, a transimpedance input stage and a
hardware discrete Fourier transform engine on a single die, so that a credible
multi-technique instrument fits on a business-card-sized board for a few tens of
dollars~\cite{devices2024ad5940, bill2023electrochemical, vamos2025hunstat2}.

What has not kept pace is the evidence offered that these instruments
\emph{measure correctly}. Reported validations lean on spectral or voltammetric
overlays, quoted error bounds, significance tests against a commercial
reference, and---for voltammetric methods above all---the coefficient of
determination of a straight-line fit against a resistive dummy
cell~\cite{pruna2018low, matsubara2021small, burgos2022tbistat,
vamos2025hunstat2}. The last of these is close to universal, because it is
cheap, requires no reference instrument, and produces a reassuring number: a
resistor obeys Ohm's law, so a correct instrument must return a straight line,
and $R^2$ values of $0.999$ and above are routinely reported as evidence of
correct operation.

This paper's central observation is that such evidence is \emph{structurally
incapable} of detecting the most consequential class of error an instrument of
this kind can have. Suppose the reported current is related to the true current
by an unknown constant gain, $I_{\mathrm{rep}} = k\,I_{\mathrm{true}}$, as it
is whenever a transimpedance gain is mis-calibrated. Against a resistor,
$I_{\mathrm{true}} = V/R$, so $I_{\mathrm{rep}} = (k/R)\,V$: still exactly a
straight line through the origin. The coefficient of determination is invariant
under any affine rescaling of the dependent variable, so $R^2$ is
\emph{identical} for every value of $k$. A linearity check cannot see a gain
error, no matter how large, because the fitted slope silently absorbs it. The
same argument disqualifies Pearson correlation, normalised RMSE, and any
"linearity (\%)" figure of merit. Only a statistic that compares the fitted
slope against an independently known quantity can detect it.

We show that this is not a hypothetical concern. Applying a traceability-first
validation protocol to a working AD5941/RP2040 instrument, we found that its
cyclic voltammetry recovered a resistance \emph{424\,\% in error}---currents low
by a factor of $5.22$---while passing a
linearity check at $R^2 > 0.99999$---and that the linearity statistic was
numerically insensitive to the error across more than two decades of induced
gain deviation. The root cause was a single software constant: the value of the
on-board calibration resistor declared in firmware did not match the resistor
actually fitted to the board. That constant sets the reference against which the
front end calibrates its own transimpedance gain, so it propagates directly into
every reported current, and it is invisible to every validation statistic in
common use.

The defect is mundane, which is precisely the point. A mismatch between a
hardware option and a firmware constant is a configuration error, not a design
error; it survives code review, unit tests, and visual inspection of the output,
because nothing about the output looks wrong. Table~\ref{tab:comparison}
summarises the validation methods reported for representative low-cost
potentiostats. Of the instruments surveyed, the voltammetric validations are
dominated by exactly the gain-blind statistics identified above.

\paragraph{Contributions} This work makes four contributions.
\begin{enumerate}
\item We give the argument, and the experimental demonstration, that linearity
  statistics are mathematically blind to transimpedance gain error, and we
  quantify the blindness by inducing gain errors of known size and showing that
  $R^2$ does not move (Section~\ref{sec:trap}).
\item We define a five-check validation protocol for voltammetric methods that
  is traceable to component nominal values and requires no reference instrument,
  and whose first check targets the scale-setting chain that linearity ignores
  (Section~\ref{sec:protocol}).
\item We apply it to the HunStat2 AD5941/RP2040 instrument on a precision
  passive Randles cell, identify and correct eight firmware defects, and report
  post-correction accuracy, repeatability, and invariance across scan rate,
  potential window and step size (Section~\ref{sec:results}).
\item We report an honest per-method coverage audit. Cyclic voltammetry is
  validated quantitatively; chronoamperometry, square-wave and differential
  pulse voltammetry, and impedance spectroscopy are shown \emph{not} to produce
  valid cell current on this hardware, with the fault localised, which
  demonstrates the companion point that one working method does not validate an
  instrument (Section~\ref{sec:coverage}).
\end{enumerate}

\begin{table*}[!t]
\caption{Validation evidence reported for representative low-cost
potentiostats. The final column asks whether the reported evidence could, even
in principle, detect a constant current-scale (transimpedance gain) error.}
\label{tab:comparison}
\centering
\footnotesize
\setlength{\tabcolsep}{4pt}
\begin{threeparttable}
\begin{tabular}{@{}l l >{\raggedright\arraybackslash}p{2.2cm} >{\raggedright\arraybackslash}p{4.3cm} c@{}}
\toprule
\textbf{Instrument} & \textbf{Front end} & \textbf{Methods} &
\textbf{Reported validation evidence} & \textbf{Gain-} \\
 & & & & \textbf{sensitive?} \\
\midrule
CheapStat~\cite{rowe2011cheapstat}      & discrete    & CV, SWV, ACV & Voltammograms of redox standards; qualitative comparison & No\tnote{a} \\
Pruna et al.~\cite{pruna2018low}        & AD5933      & EIS          & Regression against VMP3, $R^{2}=0.98$ & No\tnote{b} \\
ABE-STAT~\cite{jenkins2019abe}          & AD5933      & CV, DPV, EIS & RMS noise; potentiometric regression & No\tnote{b} \\
Matsubara~\cite{matsubara2021small}     & discrete    & CV, SWV, CA, EIS & Quoted error bounds only & Partly\tnote{c} \\
TBISTAT~\cite{burgos2022tbistat}        & AD5933      & EIS          & $t$-test; Levene; normality checks & No\tnote{d} \\
ACEStat~\cite{brown2022acestat}         & discrete    & CV, SWV, CA, EIS & Randles fitting; margin of error & Yes \\
HELPStat~\cite{bautista2025helpstat}    & AD5941      & EIS          & Paired $t$-test; Theil--Sen against reference & Yes\tnote{d} \\
HunStat2~\cite{vamos2025hunstat2}       & AD5941      & CV, OCP, EIS & Visual overlay only & No \\
\textbf{This work}                      & AD5941      & CV\tnote{e}  & Scale traceability; absolute accuracy vs.\ nominal; residual structure vs.\ LSB; replicate \%RSD; parameter invariance & Yes \\
\bottomrule
\end{tabular}
\begin{tablenotes}\footnotesize
\item[a] Qualitative waveform comparison does not constrain absolute current scale.
\item[b] $R^{2}$ and correlation are invariant under affine rescaling of the measured variable, so a constant gain error leaves them unchanged.
\item[c] Quoted bounds are gain-sensitive only if they are stated against an independently known absolute value; that basis is not reported.
\item[d] A comparison against a calibrated reference instrument is gain-sensitive in principle; a significance test alone reports whether a difference exists, not whether it is tolerable.
\item[e] CA, SWV, DPV, EIS and OCP are implemented in firmware; their status is reported in Section~\ref{sec:coverage} and they are explicitly \emph{not} claimed as validated.
\end{tablenotes}
\end{threeparttable}
\end{table*}
